\documentclass[12pt,a4paper]{article}
\usepackage[utf8]{inputenc}
\usepackage[T2A]{fontenc}
\usepackage[russian]{babel}
\usepackage{geometry}
\usepackage{float}
\geometry{top=2cm, left=2cm, right=2cm, bottom=2cm}

\title{Теория лабораторной работы 4. Метод k-ближайших соседей:\\вопросы защиты и ответы}
\author{Ануфриев А.С. 465029}
\date{07.11.2026}

\begin{document}
\maketitle

\section{Теоретические сведения}
Метод k-ближайших соседей (k-NN) --- один из простейших методов надзоруемой классификации. Объект относится к классу, который наиболее часто встречается среди k ближайших соседей в пространстве признаков. Мера близости обычно --- евклидово расстояние:
\begin{equation}
d(\mathbf{x}_i, \mathbf{x}_j) = \sqrt{\sum_{p=1}^{P} (x_{ip} - x_{jp})^2}.
\end{equation}
Параметр k определяет соседей, участвующих в голосовании: малое k делает модель чувствительной к шуму, большое k --- более устойчивой, но может сглаживать границы классов. k-NN не требует обучения в классическом смысле (ленивый алгоритм): все вычисления выполняются на этапе предсказания. Перед применением признаки масштабируются (стандартизация), иначе признаки с большими значениями доминируют в расстоянии. Качество оценивается точностью (accuracy) и матрицей ошибок.

\section{Вопросы защиты и ответы}

\subsection{Что такое метод k-ближайших соседей и к какому классу задач он относится?}
k-NN --- метод надзоруемой классификации (и регрессии), опирающийся на расстояния в пространстве признаков. Объект классифицируется по классу большинства среди k ближайших обучающих примеров.

\subsection{Как выбирается параметр k?}
k выбирается экспериментально (кросc-валидация). Малое k --- чувствительность к шуму и выбросам, большое k --- размытие границ классов. В работе рассматривались k = 3, 5, 10, 15.

\subsection{Какая метрика расстояния использовалась и почему?}
Евклидово расстояние --- стандартная метрика для k-NN на числовых признаках. Оно адекватно измеряет близость после стандартизации признаков.

\subsection{Зачем нужна стандартизация признаков перед k-NN?}
Признаки имеют разные размерности и диапазоны (например, Alcohol ~13, Proline ~747). Без масштабирования признаки с большими абсолютными значениями доминируют в расстоянии, искажая результат.

\subsection{Чем отличаются Модель 1 и Модель 2?}
Модель 1 использует случайный набор признаков (5 из 13), Модель 2 --- фиксированный набор, выбранный заранее (6 признаков). Это позволяет оценить влияние состава признаков на качество классификации.

\subsection{Что показывает матрица ошибок?}
Матрица ошибок распределяет предсказания по истинным и предсказанным классам. Диагональ --- верные предсказания, внедиагональные элементы --- ошибки классификации между конкретными классами.

\subsection{Почему ошибки классификации в работе приходятся mainly на классы 2 и 3?}
Классы 2 и 3 ближе друг к другу в пространстве признаков, чем класс 1, что затрудняет их раздел k-NN и приводит к путанице между ними.

\subsection{Какие библиотеки использовались при реализации k-NN?}
Только NumPy и Pandas (для линейной алгебры и работы с данными). Готовые средства классификации (sklearn и т.п.) не применялись.

\subsection{Какое качество показали модели?}
На тестовой выборке (36 записей) обе модели достигают точности до 0.9444. Модель 1: k=3,5 --- 0.9444; k=10,15 --- 0.9167. Модель 2: k=3,5 --- 0.8889; k=10,15 --- 0.9444.

\end{document}
